\subsection{Komposisi Fungsi}
\begin{definition}[Komposisi Fungsi]
  Komposisi dari fungsi $f$ dan $g$ didefinisikan sebagai berikut:
  \[ (f \circ g)(x) = f(g(x)) \]
  \[ (g \circ f)(x) = g(f(x)) \]
  \textit{(Menggabungkan 2 fungsi dengan cara memasukkan fungsi yang satu ke dalam variabel fungsi yang lain)}
\end{definition}

\begin{example}[Komposisi Fungsi]
  \hfill \\
  Misalkan $f(x) = x^2 + 3x - 2$ dan $g(x) = x - 1$:
  \begin{enumerate}
  \item Menentukan $(f \circ g)(x)$:
    \[
    \begin{aligned}
      (f \circ g)(x) &= f(g(x)) \\
      &= (x - 1)^2 + 3(x - 1) - 2 \\
      &= (x^2 - 2x + 1) + 3x - 3 - 2 \\
      &= x^2 + x - 4
    \end{aligned}
    \]
  \item Menentukan $(g \circ f)(x)$:
    \[
    \begin{aligned}
      (g \circ f)(x) &= g(f(x)) \\
      &= (x^2 + 3x - 2) - 1 \\
      &= x^2 + 3x - 3
    \end{aligned}
    \]
  \end{enumerate}
\end{example}

\subsection{Komposisi dari 3 Fungsi}
\begin{definition}[Komposisi dari 3 Fungsi]
  Sifat asosiatif berlaku pada komposisi fungsi:
  \[ (f \circ g) \circ h = f \circ (g \circ h) \]
\end{definition}

\begin{example}[Komposisi dari 3 Fungsi]
  \hfill \\
  Misalkan:
  \[ f(x) = x - 1 \]
  \[ g(x) = x^2 + 3x - 2 \]
  \[ h(x) = x + 5 \]
  
  Langkah 1: Tentukan $(f \circ g)(x)$:
  \[
  \begin{aligned}
    (f \circ g)(x) &= f(g(x)) \\
    &= (x^2 + 3x - 2) - 1 \\
    &= x^2 + 3x - 3
  \end{aligned}
  \]
  
  Langkah 2: Tentukan $((f \circ g) \circ h)(x)$:
  \[
  \begin{aligned}
    ((f \circ g) \circ h)(x) &= (f \circ g)(h(x)) \\
    &= (x + 5)^2 + 3(x + 5) - 3 \\
    &= (x^2 + 10x + 25) + 3x + 15 - 3 \\
    &= x^2 + 13x + 37
  \end{aligned}
  \]
\end{example}
